\scanpage{48}
% archive sheet number: 47
On va définir aussi
\[
 \deleted{\unsure{\delta}}\Delta_I\in A(I\sqcup I)
\]
qui \unsure{ser.|sera} tel que
\[
 \widetilde{\delta}_I(x)=x\qquad (x\in A(I))\qquad \text{i.e.}\quad \widetilde{\delta}_I=\mathrm{id}_{A(I)}.
\]
Pour \deleted{\unsure{$\delta$}} le voir, on \deleted{définit} considère $\delta:I\sqcup I\mathrel{\unsure{-|\to}}I$
d'où $\delta^{\deleted{*}}_*:A(I)\mathrel{-}A(I\sqcup I)$, on pose
\[
 \Delta_I=\delta_*(1)
\]
On aura (\deleted{\unsure{écrivant}} $\Delta$ au lieu de $\Delta_I$
\diagram{\begin{tikzcd}[ampersand replacement=\&,column sep=small,row sep=small]
 \&X^I\arrow[d,"\delta"]\&\\
 \&X^{I\sqcup I}\arrow[dl,"p_1"']\arrow[dr,"p_2"]\&\\
 X^I\&\&X^I
\end{tikzcd}}
\[
\begin{aligned}
 \widetilde{\deleted{\unsure{\delta}}\Delta}(x)&\overset{\text{déf de }\sim}{=}p_{2*}\bigl(\Delta\,p_1^*(x)\bigr)\\
 &=p_{2*}\bigl(\underbrace{\delta_*(1)\,p_1^*(x)}_{\parallel}\bigr)\\
 &\hspace{4em}\deleted{\unsure{p_1|p}}\delta_*\bigl(\unsure{\Gamma|\delta}^*p_1^*(x)\bigr)=\delta_*(x)\\
 &=p_{2*}\,\delta_*(x)=x
\end{aligned}
\]
Il faut encore définir l'associativité de
l'opération $\circ$, et le fait que les $\delta_I$ jouent le
rôle d'unités. On voit que sous ces conditions,

